\chapter{Quasi-isometry invariance of the Gromov boundary}
\label{chap:boundary-invariance}

Milestone M10. A quasi-isometric embedding of geodesic hyperbolic spaces induces a uniform embedding
of their Gromov boundaries (Chapter~\ref{chap:gromov-boundary}), and a quasi-isometry a uniform
equivalence (Bridson--Haefliger, Theorem III.H.3.9; Dru\c{t}u--Kapovich, Theorem 11.108;
L\"oh, Proposition 8.3.3). Isometric maps with quasi-dense range induce uniform equivalences without
any geodesic hypothesis; with the geometric realisation of Cayley graphs
(Chapter~\ref{chap:stability}) this shows that the boundary of a hyperbolic group does not depend
on the finite generating family (L\"oh, Definition 8.3.8).

The geodesic hypothesis cannot be dropped for quasi-isometries: Bridson--Haefliger
(Example III.H.3.14) give a hyperbolic subset of the plane, quasi-isometric to a ray, whose boundary
has two points while that of the ray has one.

The proof follows Dru\c{t}u--Kapovich's proof of Theorem 11.108: a quasi-isometric embedding of
geodesic spaces with slim triangles distorts Gromov products at most affinely, in both directions,
because the Gromov product $(x \cdot y)_p$ is, up to $2\delta$, the distance from $p$ to a geodesic
from $x$ to $y$, and geodesics are carried to quasi-geodesics close to geodesics (stability,
Theorem~\ref{Metric.HasSlimTrianglesWith.exists_hausdorffDist_le}). Working with the uniform
structure of the boundary avoids their use of sequences, properness and compactness.

\paragraph{Differences from the proofs with rays.} L\"oh and Bridson--Haefliger define the boundary
by rays and prove the invariance there.
\begin{itemize}
  \item \emph{L\"oh} (Proposition 8.3.3, boundary of quasi-geodesic rays up to finite Hausdorff
        distance). The map $[\gamma] \mapsto [f \circ \gamma]$ is well defined, injective, and
        unchanged when $f$ moves a bounded distance, by the bounds of $f$ alone, with no
        hyperbolicity; so a quasi-isometry induces a bijection. Continuity is left as an exercise
        (8.E.14), and it is where the work lies: the topology is uniform convergence of
        representatives on compact sets, and $f$ is not continuous, so $f \circ \gamma_n$ need not
        converge to $f \circ \gamma$ when $\gamma_n$ converges to $\gamma$; other representatives,
        close to geodesic rays, have to be found.
  \item \emph{Bridson--Haefliger} (Theorem III.H.3.9, geodesic rays in proper geodesic hyperbolic
        spaces). Each quasi-geodesic ray is first replaced by a geodesic ray at finite Hausdorff
        distance, starting at any given point (Lemma III.H.3.1: geodesics from the base point to
        points far along the ray stay close to it by stability, and Arzel\`a--Ascoli, which needs
        properness, gives a limiting ray). The class of a ray $c$ is sent to that of the geodesic
        ray straightening $f \circ c$. Continuity uses the neighbourhoods of Lemma III.H.3.6 (rays
        from the base point within a fixed distance of a given ray at time $n$): rays close up to
        time $n$ have images close up to time about $n/\lambda$, and stability with thin
        triangles makes their straightenings fellow-travel up to a time that tends to infinity
        with $n$.
\end{itemize}
Both routes use stability, as the route here does. The proof here does not need the straightening
of rays, and so neither properness nor Arzel\`a--Ascoli, nor a topology built from convergent
sequences of rays: the estimate of
Lemma~\ref{Metric.IsQuasiIsometricEmbedding.exists_gromovProduct_le} gives uniform continuity
directly. Dru\c{t}u--Kapovich also prove their Theorem 11.108 on the sequential boundary and
transfer it to rays through the homeomorphism of the two boundaries (their Theorem 11.104).

How boundary maps are defined is an open design decision (D32); the nodes below follow its
proposal and are not implemented until it is decided.

\section*{Gromov products and geodesics}

In any geodesic space, $(x \cdot y)_p \le d(p, q)$ for every point $q$ of a geodesic from $x$ to
$y$, by the triangle inequality. With slim triangles there is a converse:

\begin{lemma}[cf.\ Dru\c{t}u--Kapovich Lem.~11.22]
  \label{Metric.HasSlimTrianglesWith.infDist_le_gromovProduct_add}
  \uses{Metric.HasSlimTrianglesWith, Metric.Geodesic}
  Let $X$ be a geodesic pseudometric space with $\delta$-slim triangles, $\gamma$ a geodesic from
  $x$ to $y$ and $p \in X$. Some point of $\gamma$ is within $(x \cdot y)_p + 2\delta$ of $p$.
\end{lemma}
\begin{proof}
  Choose geodesics from $y$ to $p$ and from $p$ to $x$; with $\gamma$ they form a $\delta$-slim
  triangle. The parameters $t$ for which $\gamma(t)$ is within $\delta$ of the side from $p$ to $x$,
  respectively from $y$ to $p$, form two closed sets covering the interval of $\gamma$ and
  containing its left, respectively right, endpoint; since the interval is connected, some $z =
  \gamma(t)$ is within $\delta$ of a point $a$ of the first side and a point $b$ of the second. Then
  $d(p, z) \le d(p, a) + \delta = d(p, x) - d(a, x) + \delta \le d(p, x) - d(x, z) + 2\delta$, and
  likewise $d(p, z) \le d(p, y) - d(y, z) + 2\delta$. Adding, with $d(x, z) + d(z, y) = d(x, y)$,
  gives $d(p, z) \le (x \cdot y)_p + 2\delta$. This is the step (b) of the proof of
  Lemma~\ref{Metric.HasSlimTrianglesWith.isHyperbolicWith}.
\end{proof}

\begin{lemma}[Quasi-isometric embeddings and Gromov products]
  \label{Metric.IsQuasiIsometricEmbedding.exists_gromovProduct_le}
  \uses{Metric.IsQuasiIsometricEmbedding, Metric.Geodesic, Metric.HasSlimTrianglesWith}
  Let $X$ and $Y$ be geodesic pseudometric spaces, $Y$ with $\delta$-slim triangles, and $f \colon
  X \to Y$ a quasi-isometric embedding. There are $A, B \ge 0$ such that for all $p, x, y \in X$
  \[ (x \cdot y)_p \le A\,(f x \cdot f y)_{f p} + B \quad\text{and}\quad
     (f x \cdot f y)_{f p} \le A\,(x \cdot y)_p + B. \]
  \plannedlean{Metric.IsQuasiIsometricEmbedding.exists_gromovProduct_comp_le}
\end{lemma}
\begin{proof}
  \uses{Metric.HasSlimTrianglesWith.infDist_le_gromovProduct_add,
    Metric.HasSlimTrianglesWith.exists_hausdorffDist_le, Metric.QuasiGeodesic,
    Metric.IsQuasiIsometricEmbedding.exists_hasSlimTrianglesWith}
  Let $f$ be a $(c, b)$-quasi-isometric embedding, so $d(f u, f v) \le c\,d(u, v) + b$ and $d(u, v)
  \le c\,d(f u, f v) + cb$. For a geodesic $\gamma$ from $x$ to $y$, $f \circ \gamma$ is a $(c,
  b)$-quasi-geodesic from $f x$ to $f y$; let $R$ be the stability constant of $Y$ for $c, b$
  (Theorem~\ref{Metric.HasSlimTrianglesWith.exists_hausdorffDist_le}) and $\gamma'$ a geodesic from
  $f x$ to $f y$. Every point of $\gamma'$ is within $R + 1$ of a point of $f \circ \gamma$, and
  conversely.

  \emph{First inequality.} Lemma~\ref{Metric.HasSlimTrianglesWith.infDist_le_gromovProduct_add}
  gives $q'$ on $\gamma'$ with $d(f p, q') \le (f x \cdot f y)_{f p} + 2\delta$; take $q$ on
  $\gamma$ with $d(f q, q') \le R + 1$. Then $(x \cdot y)_p \le d(p, q) \le c\,d(f p, f q) + cb \le
  c\,(f x \cdot f y)_{f p} + c(2\delta + R + 1 + b)$.

  \emph{Second inequality.} $X$ has $\delta'$-slim triangles for some $\delta'$
  (Corollary~\ref{Metric.IsQuasiIsometricEmbedding.exists_hasSlimTrianglesWith}), so there is $q$ on
  $\gamma$ with $d(p, q) \le (x \cdot y)_p + 2\delta'$; take $q'$ on $\gamma'$ with $d(f q, q') \le
  R + 1$. Then $(f x \cdot f y)_{f p} \le d(f p, q') \le c\,d(p, q) + b + R + 1 \le c\,(x \cdot y)_p
  + 2c\delta' + b + R + 1$.

  Take $A = \max\{c, 1\}$ and $B$ the larger additive constant.
\end{proof}

\section*{Boundary maps}

\begin{definition}[Boundary map]
  \label{Metric.GromovBoundary.map}
  \uses{Metric.GromovBoundary, Metric.EquivAtInfinity}
  Let $X$ and $Y$ be hyperbolic and $f \colon X \to Y$ a map that sends sequences equivalent at
  infinity to sequences equivalent at infinity. Then $f$ sends sequences converging at infinity to
  sequences converging at infinity, and $\partial f \colon \partial X \to \partial Y$, $[x_n]
  \mapsto
  [f(x_n)]$, is well defined.
\end{definition}

\begin{lemma}[Functoriality, cf.\ L\"oh Prop.~8.3.3 (2)]
  \label{Metric.GromovBoundary.map_eq_of_hasFiniteDist}
  \uses{Metric.GromovBoundary.map, Metric.HasFiniteDist}
  Maps $f, g \colon X \to Y$ as in Definition~\ref{Metric.GromovBoundary.map} that have finite
  distance induce the same boundary map. Moreover
  $\partial(g \circ f) = \partial g \circ \partial f$.
  \plannedlean{Metric.GromovBoundary.map_comp}
\end{lemma}
\begin{proof}
  Sequences $(u_n)$, $(v_n)$ with $d(u_n, v_n) \le C$ and $(u_n)$ converging at infinity are
  equivalent at infinity, as $(u_i \cdot v_j)_q \ge (u_i \cdot u_j)_q - C$. Apply this to $u_n =
  f(x_n)$, $v_n = g(x_n)$. The composition holds by definition.
\end{proof}

\begin{theorem}[Extension theorem, Bridson--Haefliger Thm.~III.H.3.9, Dru\c{t}u--Kapovich
  Thm.~11.108, L\"oh Prop.~8.3.3 (1)]
  \label{Metric.IsQuasiIsometricEmbedding.isUniformEmbedding_gromovBoundaryMap}
  \uses{Metric.GromovBoundary.map, Metric.GromovBoundary.instUniformSpace,
    Metric.IsQuasiIsometricEmbedding, Metric.Geodesic}
  Let $X$ and $Y$ be geodesic hyperbolic pseudometric spaces and $f \colon X \to Y$ a
  quasi-isometric embedding. Then $f$ sends sequences to sequences equivalent at infinity exactly
  when they are, and $\partial f$ is a uniform embedding.
  \plannedlean{Metric.IsQuasiIsometricEmbedding.equivAtInfinity_comp_iff}
\end{theorem}
\begin{proof}
  \uses{Metric.IsQuasiIsometricEmbedding.exists_gromovProduct_le,
    Metric.IsHyperbolicWith.hasSlimTrianglesWith, Metric.GromovBoundary.min_gromovProduct_le}
  $Y$ has slim triangles (Lemma~\ref{Metric.IsHyperbolicWith.hasSlimTrianglesWith}); let $A, B$ be
  as in Lemma~\ref{Metric.IsQuasiIsometricEmbedding.exists_gromovProduct_le}, and $\delta$ a
  four-point constant of $Y$. The two inequalities, at the base points $p$ and $f p$, show that
  $(x_i \cdot y_j)_p \to \infty$ if and only if $(f x_i \cdot f y_j)_{f p} \to \infty$; so $f$
  preserves and reflects equivalence at infinity, and $\partial f$ is defined and injective. For
  representatives $(x_n)$ of $\xi$ and $(y_n)$ of $\eta$, the sequences $(f x_n)$, $(f y_n)$
  represent $\partial f \xi$, $\partial f \eta$, so
  \[ (\partial f \xi \cdot \partial f \eta)_{f p} \ge \liminf (f x_i \cdot f y_j)_{f p} \ge
     \bigl(\liminf (x_i \cdot y_j)_p - B\bigr)/A; \]
  taking the supremum over representatives, $(\partial f \xi \cdot \partial f \eta)_{f p} \ge ((\xi
  \cdot \eta)_p - B)/A$. Conversely, by
  Lemma~\ref{Metric.GromovBoundary.min_gromovProduct_le} (2) in $Y$,
  $(\partial f \xi \cdot \partial f \eta)_{f p} \le \liminf (f x_i \cdot f y_j)_{f p} + 2\delta \le
  A\,(\xi \cdot \eta)_p + B + 2\delta$. So $\partial f$ maps $V_{p, AR + B}$ into $V_{f p, R}$, and
  the preimage of $V_{f p, A R + B + 2\delta}$ lies in $V_{p, R}$; as all base points give the same
  uniform structure, $\partial f$ is uniformly continuous and induces the uniform structure of
  $\partial X$.
\end{proof}

\begin{corollary}[Quasi-isometries, Bridson--Haefliger Thm.~III.H.3.9, L\"oh Prop.~8.3.3]
  \label{Metric.IsQuasiIsometry.gromovBoundaryEquiv}
  \uses{Metric.IsQuasiIsometry, Metric.GromovBoundary.map, Metric.GromovBoundary.instUniformSpace,
    Metric.Geodesic}
  A quasi-isometry $f \colon X \to Y$ of geodesic hyperbolic pseudometric spaces induces a uniform
  equivalence $\partial f \colon \partial X \to \partial Y$, in particular a homeomorphism.
\end{corollary}
\begin{proof}
  \uses{Metric.IsQuasiIsometricEmbedding.isUniformEmbedding_gromovBoundaryMap,
    Metric.GromovBoundary.map_eq_of_hasFiniteDist}
  Let $g$ be a quasi-inverse, also a quasi-isometric embedding. By
  Lemma~\ref{Metric.GromovBoundary.map_eq_of_hasFiniteDist}, $\partial g \circ \partial f =
  \partial(g \circ f) = \partial \mathrm{id} = \mathrm{id}$ and likewise $\partial f \circ \partial
  g = \mathrm{id}$; both maps are uniformly continuous by
  Theorem~\ref{Metric.IsQuasiIsometricEmbedding.isUniformEmbedding_gromovBoundaryMap}.
\end{proof}

\begin{corollary}[Isometric maps with quasi-dense range]
  \label{Isometry.gromovBoundaryEquiv}
  \uses{Metric.GromovBoundary.map, Metric.GromovBoundary.instUniformSpace, Metric.QuasiDenseRange}
  Let $X$ and $Y$ be hyperbolic pseudometric spaces and $f \colon X \to Y$ an isometric map with
  quasi-dense range. Then $\partial f$ is a uniform equivalence.
\end{corollary}
\begin{proof}
  \uses{Metric.GromovBoundary.min_gromovProduct_le}
  $f$ preserves Gromov products, hence equivalence at infinity in both directions; so $\partial f$
  is defined and injective, and $(\xi \cdot \eta)_p \le (\partial f \xi \cdot \partial f \eta)_{f p}
  \le (\xi \cdot \eta)_p + 2\delta$ as in the proof of
  Theorem~\ref{Metric.IsQuasiIsometricEmbedding.isUniformEmbedding_gromovBoundaryMap}, so it is a
  uniform embedding. Surjective: let every point of $Y$ be within $C$ of the range and $(y_n)$
  converge at infinity; choose $x_n$ with $d(f x_n, y_n) \le C$. Then $(f x_n)$ is equivalent to
  $(y_n)$ (as in the proof of Lemma~\ref{Metric.GromovBoundary.map_eq_of_hasFiniteDist}), so it
  converges at infinity, and so does $(x_n)$; and $\partial f [x_n] = [y_n]$.
\end{proof}

\section*{Hyperbolic groups}

L\"oh (Definition 8.3.8) defines the boundary of a finitely generated group as that of its Cayley
graph for a finite generating set. Here it is the boundary of $(G, d_P)$, the group with the word
metric of a finite generating family $P$; no new definition is needed. It is the boundary of the
geometric realisation of the Cayley graph, by Corollary~\ref{Isometry.gromovBoundaryEquiv} applied
to
the vertex map.

\begin{theorem}[Independence of the generating family, cf.\ L\"oh Def.~8.3.8]
  \label{Group.Generators.WithWordMetric.gromovBoundaryEquiv}
  \uses{Group.Generators.WithWordMetric, Group.Generators.IsHyperbolicWith, Metric.GromovBoundary,
    Metric.GromovBoundary.instUniformSpace}
  Let $G$ be a hyperbolic group. For every finite generating family $P$ the word metric of $P$ is
  hyperbolic, and for finite generating families $P$, $Q$ there is a uniform equivalence
  $\partial(G, d_P) \to \partial(G, d_Q)$ sending $[x_n]$ to $[x_n]$.
  \plannedlean{Group.Generators.WithWordMetric.instIsHyperbolic}
\end{theorem}
\begin{proof}
  \uses{Group.Generators.isHyperbolic_iff, SimpleGraph.Realization,
    SimpleGraph.Realization.dist_vertex, SimpleGraph.Realization.isGeodesicSpace,
    Group.Generators.dist_cayleySimpleGraph, Metric.IsHyperbolicWith.of_isometry_of_quasiDenseRange,
    Isometry.gromovBoundaryEquiv, Metric.IsQuasiIsometry.gromovBoundaryEquiv,
    Group.Generators.WithWordMetric.isQuasiIsometry_congr,
    Metric.isQuasiIsometry_iff_quasiDenseRange,
    Metric.IsQuasiIsometry.comp, Metric.GromovBoundary.map_eq_of_hasFiniteDist}
  The word metric of $P$ is hyperbolic by Theorem~\ref{Group.Generators.isHyperbolic_iff}. Let $|P|$
  be the geometric realisation of the Cayley graph of $P$ and $\iota_P \colon (G, d_P) \to |P|$ the
  vertex map: isometric with quasi-dense range
  (Lemma~\ref{SimpleGraph.Realization.dist_vertex},
  Theorem~\ref{Group.Generators.dist_cayleySimpleGraph}), so $|P|$ is hyperbolic
  (Lemma~\ref{Metric.IsHyperbolicWith.of_isometry_of_quasiDenseRange}) and geodesic
  (Lemma~\ref{SimpleGraph.Realization.isGeodesicSpace}), and $\partial \iota_P$ is a uniform
  equivalence (Corollary~\ref{Isometry.gromovBoundaryEquiv}). The same holds for $Q$. Let $h_P$ be a
  quasi-inverse of $\iota_P$ with $h_P \circ \iota_P = \mathrm{id}$. Then $\iota_Q \circ \mathrm{id}
  \circ h_P \colon |P| \to |Q|$ is a quasi-isometry
  (Corollary~\ref{Group.Generators.WithWordMetric.isQuasiIsometry_congr},
  Lemma~\ref{Metric.IsQuasiIsometry.comp}) of geodesic hyperbolic spaces, so it induces a uniform
  equivalence (Corollary~\ref{Metric.IsQuasiIsometry.gromovBoundaryEquiv}). The composite
  $(\partial \iota_Q)^{-1} \circ \partial(\iota_Q \circ h_P) \circ \partial \iota_P$ sends $[x_n]$
  to $[x_n]$, since $h_P \circ \iota_P = \mathrm{id}$ and boundary maps compose
  (Lemma~\ref{Metric.GromovBoundary.map_eq_of_hasFiniteDist}).
\end{proof}
